\DocumentMetadata{
  lang        = en-US,
  tagging     = on,
}

\RequirePackage{shellesc}
\documentclass[11pt, oneside]{article}   	% use "amsart" instead of "article" for AMSLaTeX format

\usepackage{geometry}                		% See geometry.pdf to learn the layout options. There are lots.
\geometry{letterpaper}                   		% ... or a4paper or a5paper or ... 
%\geometry{landscape}                		% Activate for rotated page geometry
\usepackage[parfill]{parskip}    		% Activate to begin paragraphs with an empty line rather than an indent
\usepackage{graphicx}				% Use pdf, png, jpg, or eps§ with pdflatex; use eps in DVI mode
								% TeX will automatically convert eps --> pdf in pdflatex		
\usepackage{amssymb}
\usepackage{amsmath}
\usepackage{tablefootnote}
\usepackage{multirow}
\usepackage{tabularx}
\renewcommand\tabularxcolumn[1]{m{#1}}

\usepackage{dsfont}
\usepackage[hang,flushmargin]{footmisc}
\usepackage[normalem]{ulem}

\usepackage{pgfplots,pgfplotstable}

\usepackage[autosize,forceshell]{dot2texi}
\usepackage{tikz}
\usetikzlibrary{shapes,arrows,patterns}
\usepackage{tkz-fct}

\ExplSyntaxOn
\cs_new_protected:Npn \suspendtagging #1 { \tag_suspend:n{#1} }
\cs_new_protected:Npn \resumetagging  #1 { \tag_resume:n{#1} }
\ExplSyntaxOff

\makeatletter
\long\def\ifnodedefined#1#2#3{%
    \@ifundefined{pgf@sh@ns@#1}{#3}{#2}%
}

\pgfplotsset{
    discontinuous/.style={
    scatter,
    scatter/@pre marker code/.code={
        \ifnodedefined{marker}{
            \pgfpointdiff{\pgfpointanchor{marker}{center}}%
             {\pgfpoint{0}{0}}%
             \ifdim\pgf@y>0pt
                \tikzset{options/.style={mark=*}}
                \draw [dashed,thin] (marker-|0,0) -- (0,0);
                \draw plot [mark=*,mark options={fill=white}] coordinates {(marker-|0,0)};
             \else
                \tikzset{options/.style={mark=none}}
             \fi
        }{
            \tikzset{options/.style={mark=none}}        
        }
        \coordinate (marker) at (0,0);
        \begin{scope}[options]
    },
    scatter/@post marker code/.code={\end{scope}}
    }
}

\makeatother

\newcommand*\Eval[3]{\left(#1\middle)\right\rvert_{#2}^{#3}}


\usepackage{arydshln}
\usepackage{mathtools}


%\usepackage{tagpdf}
%\tagpdfsetup{activate,paratagging,interwordspace}
\usepackage{accsupp}
\usepackage{hyperref}
\usepackage{lipsum}
%\usepackage{axessibility}
%\usepackage{accessibility}

%SetFonts
\DeclareEmphSequence{\bfseries\itshape}
\DeclareMathOperator {\ints}{\mathbb{Z}}
\DeclareMathOperator {\reals}{\mathbb{R}}
\DeclareMathOperator*{\argmin}{argmin}
\DeclareMathOperator*{\argmax}{argmax}
\DeclareMathOperator{\E}{\mathbb{E}}
\DeclareMathOperator{\V}{\mathbb{V}}


\renewcommand\labelenumi{(\theenumi)}
\renewcommand*{\thefootnote}{\fnsymbol{footnote}}

\usepackage{tcolorbox}

\newenvironment{aside}[1][Aside]{\begin{tcolorbox}[colback=black!5,colframe=black!75!black,title=#1]}{\end{tcolorbox}}

%SetFonts


\title{Topic 13: Introduction to Probability}
\author{02-680: Essentials of Mathematics and Statistics}
%\date{}							% Activate to display a given date or no date

\begin{document}
\maketitle

\begin{center}
Probability is the mathematical language for quantifying uncertainty
\end{center} 

%%%%%%%%%%%%
\section{Basics}
The \emph{sample space} ($\Omega$), sometimes called the universe, is the set of all possible
outcomes of an experiment (or random process).
The outcomes are \textit{mutually exclusive}: exactly one outcome happens each time,
and together they cover everything that could happen.
(This $\Omega$ is not the $\Omega$ from asymptotics; same letter, different object.)
\[\Omega_{coin}=\{\texttt{Heads},\texttt{Tails}\}\]
\[\Omega_{twocoin}=\Omega_{coin}\times\Omega_{coin}=\{
\langle\texttt{Heads},\texttt{Heads}\rangle,
\langle\texttt{Heads},\texttt{Tails}\rangle,
\langle\texttt{Tails},\texttt{Heads}\rangle,
\langle\texttt{Tails},\texttt{Tails}\rangle\}\]

A sample \emph{outcome}, or atomic element, $\omega\in\Omega$ is one of the possible things that can happen. 

An \emph{event} is a \emph{subset} of the sample space: a group of outcomes that we bundle together because we care about them as a group. 
Outcomes are the \emph{elements} of $\Omega$; events are the \emph{subsets} of $\Omega$. 
For instance, the event that the first of two coin tosses is heads: 
\[A_{firstheads}=\{
\langle\texttt{Heads},\texttt{Heads}\rangle,
\langle\texttt{Heads},\texttt{Tails}\rangle\} \subseteq \Omega_{twocoin} \]

We say the \emph{probability} $p(A)$ of an event $A$ is the fraction of all outcomes that $A$ covers.
(For now we assume all outcomes are equally likely; the axioms below describe what $p$ must satisfy in general.)
So in the example above $p(A_{firstheads})=\frac{1}{2}$.
Notice that $p$ takes an \emph{event} (a set), not an outcome. 
When we write $p(\texttt{Heads})$ we really mean $p(\{\texttt{Heads}\})$, the event containing only that outcome. 

\paragraph{Event space.}
The \emph{event space}, $\mathcal{A}$, is the set of all events we are allowed to ask about.
That makes it a set of sets.
(We use the calligraphic $\mathcal{A}$ for the set of events and a plain $A$ for a single event.)
For a single coin, $\Omega_{coin}=\{\texttt{Heads},\texttt{Tails}\}$ has two outcomes, but four events:
\[\mathcal{A}_{coin}=\Big\{\;\emptyset,\;\;\{\texttt{Heads}\},\;\;\{\texttt{Tails}\},\;\;\{\texttt{Heads},\texttt{Tails}\}\;\Big\}\]
(``nothing happens,'' ``heads,'' ``tails,'' and ``heads or tails''). 
For discrete sample spaces (such as coin flips) we can typically take $\mathcal{A}$ to be the power set of $\Omega$, 
so $|\mathcal{A}|=2^{|\Omega|}$. 
The two-coin sample space has $|\Omega_{twocoin}|=4$ outcomes and so $2^4=16$ events. 
In continuous spaces it is typically thought of as the Borel field of $\Omega$ (details of this are beyond the scope of the class for now). 

\begin{center}
\begin{tabular}{|l|l|l|}
\hline
 & sample space $\Omega$ & event space $\mathcal{A}$\\\hline
 contains & outcomes (what can happen) & events (what we can ask about)\\\hline
 its members are & elements $\omega\in\Omega$ & subsets $A\subseteq\Omega$, so $A\in\mathcal{A}$\\\hline
 one coin & 2 outcomes & 4 events\\\hline
 two coins & 4 outcomes & 16 events\\\hline
\end{tabular}
\end{center}

\begin{aside}[Think of it this way]
$\Omega$ is the list of everything that could be in the freezer: each tissue sample is an outcome. 
$\mathcal{A}$ is the list of every yes/no question you could ask about the sample you pulled, such as ``is it high expression?'' 
Each question is the \emph{set} of outcomes for which the answer is yes, so each question is an event. 
\end{aside}

\paragraph{Why the event space has rules.}
If we can ask whether $A$ happened, we should also be able to ask whether it did \emph{not} happen, and whether $A$ or $B$ happened:
\begin{enumerate}
\item $A\in\mathcal{A} \rightarrow \overline{A}\in\mathcal{A}$ \quad (we can ask ``not $A$''; the \emph{complement} $\overline{A}=\Omega\setminus A$ is everything in $\Omega$ that is not in $A$)
\item $A_1,A_2,... \in \mathcal{A} \rightarrow \bigcup_{i=1}^\infty A_i \in \mathcal{A}$ \quad (we can ask ``$A_1$ or $A_2$ or ...'')
\end{enumerate}
These are exactly what the axioms below need: Axiom 3 adds probabilities of unions, 
and the complement rule $p(\overline{A})=1-p(A)$ only makes sense if $\overline{A}$ is an event. 

%%%%%%%%
\subsection{Three Axioms}

\paragraph{Axiom 1: Non-negative Probability.}
For any event $A$, $p(A)\ge 0$.

That is, we cannot have a negative probability (zero is allowed).
In the previous example if $A_{threetails}$ is the event that you get 3 tails when a coin is tossed twice,
then $A_{threetails}=\emptyset$ (it can't happen), so $p(A_{threetails})=0$ but this can't be negative.

\paragraph{Axiom 2: Total Probability.}
For any sample space $\Omega$, $p(\Omega)=1$. 
That is, the probability of something in the sample space happening is 1. 
So for the example above the probability of a single coin flip being either \texttt{Heads} or \texttt{Tails} is 1, 
(as defined) there are no other possible outcomes. 

\paragraph{Axiom 3: Disjoint Events.}
For \textit{pairwise disjoint} events $A_1, A_2, ..., A_n$ (no two of them share an outcome), \[p(A_1\cup A_2\cup ... \cup A_n)=p(A_1) + p(A_2) + ... + p(A_n).\]
So in a single coin flip if $A_{heads}$ and $A_{tails}$ are the events of a single coin flip being heads and tails respectively, 
we can see the events are disjoint (they don't share any outcomes)
so \[p(A_{heads} \cup A_{tails}) = p(A_{heads}) + p(A_{tails}) = \frac{1}{2} + \frac{1}{2} = 1.\]

As a counter example lets define \[A_{lastheads}=\{
\langle\texttt{Heads},\texttt{Heads}\rangle,
\langle\texttt{Tails},\texttt{Heads}\rangle\} \subset \Omega_{twocoin}. \]
We can see that 
\[p(A_{firstheads} \cup A_{lastheads}) \ne p(A_{firstheads}) + p(A_{lastheads})\]
because $A_{firstheads} \cap A_{lastheads} = \{\langle\texttt{Heads},\texttt{Heads}\rangle\} \ne \emptyset$ thus they are not disjoint.
In fact the union has only 3 outcomes, so $p(A_{firstheads} \cup A_{lastheads})=\frac{3}{4}$,
but $p(A_{firstheads}) + p(A_{lastheads}) = \frac{1}{2}+\frac{1}{2} = 1$:
the outcome $\langle\texttt{Heads},\texttt{Heads}\rangle$ was counted twice.


%%%%%%%%
\subsection{Example Consequences of the Three Axioms}

\subsubsection{$p(\emptyset) = 0$}
$\Omega \cup \emptyset = \Omega$. 

Since $\Omega$ is disjoint from $\emptyset$, 
we know $p(\Omega\cup\emptyset) = p(\Omega) + p(\emptyset)$. 

Put these two together we see that 
$p(\Omega) = p(\Omega) + p(\emptyset)$. 

Since $p(\Omega)=1$, 
then $1 = 1 + p(\emptyset)$, 
and thus $p(\emptyset) = 0$. 

\subsubsection{$p(\overline{A}) = 1-p(A)$}
We know by definition that $\overline{A} = \Omega \setminus A$. 

From Axiom 3, since $A$ and $\overline{A}$ are disjoint, 
we know $p(A\cup\overline{A}) = p(A) + p(\overline{A})$.

But we also know that $A\cup\overline{A} = \Omega$. 

So that means $ p(A) + p(\overline{A}) = p(\Omega) = 1$. 
With some algebra we find the original statement. 

\subsubsection{$0 \le p(A) \le 1$.}

From axiom 1: $0 \le p(A)$.

For the other side, the previous result gives $p(A) = 1-p(\overline{A})$,
and $p(\overline{A})\ge 0$ by axiom 1, so $p(A) \le 1-0 = 1$.

\subsubsection{If $A\subseteq B$ then $p(A) \le p(B)$.}

\begin{align*}
p(B) & = p(A \cup (B\setminus A)) &			B = A \cup (B\setminus A)\\
	&= p(A) + p(B\setminus A) & A \textnormal{ and } B\setminus A \textnormal{ are disjoint (Axiom 3)} \\
	&\ge p(A) 				& p(B\setminus A) \ge 0 \textnormal{ (Axiom 1)} \\
\end{align*}

%%%%%%%%%%%%%
\section{Perspectives}

For a coin toss, there are two common ways of explaining what $p(\texttt{Heads})=\frac{1}{2}$ means:

\paragraph{Frequentists.}
Probability is a long-run frequency.
You flip a coin 100 times, you get \texttt{Heads} approximately 50 of those tosses,
and the more you flip the closer the fraction of \texttt{Heads} gets to $\frac{1}{2}$.

\paragraph{Bayesian.}
Probability is a degree of belief.
You \textit{believe} you will get heads 50\% of the times you flip the coin,
and you update that belief as you see evidence (for instance, a coin that keeps landing \texttt{Tails}).

Both views follow the same three axioms and the same rules derived from them,
so everything in this topic holds either way; they differ in what the number $p(A)$ is said to \emph{mean}.

\section*{Useful References}
Wasserman. ``All of Statistics: A Concise Course in Statistical Inference'' \S\S 1.2--1.3\\
Degroot and Schervish. ``Probability and Statistics'' \S\S 1.2--1.5

\end{document}
